% Options for packages loaded elsewhere
\PassOptionsToPackage{unicode}{hyperref}
\PassOptionsToPackage{hyphens}{url}
\PassOptionsToPackage{dvipsnames,svgnames,x11names}{xcolor}
%
\documentclass[
  11pt,
]{article}
\usepackage{amsmath,amssymb}
\usepackage{setspace}
\usepackage{iftex}
\ifPDFTeX
  \usepackage[T1]{fontenc}
  \usepackage[utf8]{inputenc}
  \usepackage{textcomp} % provide euro and other symbols
\else % if luatex or xetex
  \usepackage{unicode-math} % this also loads fontspec
  \defaultfontfeatures{Scale=MatchLowercase}
  \defaultfontfeatures[\rmfamily]{Ligatures=TeX,Scale=1}
\fi
\usepackage[]{newtx}
\ifPDFTeX\else
  % xetex/luatex font selection
\fi
% Use upquote if available, for straight quotes in verbatim environments
\IfFileExists{upquote.sty}{\usepackage{upquote}}{}
\IfFileExists{microtype.sty}{% use microtype if available
  \usepackage[]{microtype}
  \UseMicrotypeSet[protrusion]{basicmath} % disable protrusion for tt fonts
}{}
\makeatletter
\@ifundefined{KOMAClassName}{% if non-KOMA class
  \IfFileExists{parskip.sty}{%
    \usepackage{parskip}
  }{% else
    \setlength{\parindent}{0pt}
    \setlength{\parskip}{6pt plus 2pt minus 1pt}}
}{% if KOMA class
  \KOMAoptions{parskip=half}}
\makeatother
\usepackage{xcolor}
\usepackage[margin=1in]{geometry}
\setlength{\emergencystretch}{3em} % prevent overfull lines
\providecommand{\tightlist}{%
  \setlength{\itemsep}{0pt}\setlength{\parskip}{0pt}}
\setcounter{secnumdepth}{-\maxdimen} % remove section numbering
\ifLuaTeX
  \usepackage{selnolig}  % disable illegal ligatures
\fi
\usepackage{bookmark}
\IfFileExists{xurl.sty}{\usepackage{xurl}}{} % add URL line breaks if available
\urlstyle{same}
\hypersetup{
  colorlinks=true,
  linkcolor={Maroon},
  filecolor={Maroon},
  citecolor={Blue},
  urlcolor={Blue},
  pdfcreator={LaTeX via pandoc}}

\author{}
\date{}

\begin{document}

\setstretch{1.15}
\section{Response to the Operations Research
Referee}\label{response-to-the-operations-research-referee}

\textbf{Manuscript:} Accepted Service Adaptation: Compact Allocation
Quotients and Minimal Price Memory

\textbf{Revision:} R30, September 23, 2026

\textbf{Report answered:} Price-state R29 report of September 23, 2026.

\textbf{Review commit:}
\texttt{d1f7f977132e46ac9ca582a1d78f5219365509c8}.

\textbf{Reviewed scientific baseline:}
\texttt{51552e4198f09b82c764956b92d7b1ee2ad57aff}.

\subsection{Overview of the revision}\label{overview-of-the-revision}

We thank the referee for identifying the direct tree-allocation
connection and for distinguishing the substantive finite-price result
from the earlier nominal revision. We have rebuilt the publication
manuscript around the exact laminar mapping, the compact quotient, its
rational complexity, and minimal execution memory. The response does not
replace a mathematical objection with a change in terminology. The new
manuscript contains proofs of the additional results, implements the
coefficient-event algorithm needed by the bit bound, and compares three
explicitly specified alternatives.

The principal additions are Proposition 2.1 (exact laminar mapping),
Theorem 3.1 (occurrence-to-vertex quotient with the finite response
closure), Theorem 3.3 (polynomial rational coefficient and preprocessing
bounds), Theorem 4.1 (minimal behavioral alphabet, including randomized
exact-optimum implementation), Theorem 4.2 (heterogeneous ordered memory
frontier), and Proposition 4.3 (outside-option dispersion at a fixed
mean payment). Proposition 5.2 completes the endpoint-normal argument
for supporting cuts. All previously established mathematics and
numerical records remain available in unchanged historical sources and
byte-identical predecessor copies.

The new evidence contains 54 completed comparison workers and 12 scaling
workers. A separate final audit recompiles all 22 distinct saved
certificates from the final code. The largest completed public graph has
2,045 vertices and 8,164 edges; a separate 511-vertex family stores
197,369 response segments. This distinction prevents a low-breakpoint
large graph from being presented as a worst-case storage test.

\subsection{1. Version audit and scientific
delta}\label{version-audit-and-scientific-delta}

We agree that the scientifically relevant starting point is the
price-state R29 manuscript, not the nominal R29 branch. The source
package was reconstructed against the exact upstream tree and commit,
and the revised branch descends from that scientific baseline. The
current response addresses the subsequent price-state report, including
its classical-allocation objection. \texttt{PROVENANCE.json} records
both distinct source and review commits. The historical report and other
revision branches are not edited.

The additional results are identified by theorem rather than by the
number of files or pages changed. \texttt{NOVELTY\_MATRIX.md} separates
inherited principles, the R29 theorem, R30 extensions, and unclaimed
generalizations.

\subsection{2. What the central theorem
proves}\label{what-the-central-theorem-proves}

Theorem 3.1 now begins with the occurrence quotient itself. It defines
an occurrence, the public vertex, incoming price, relative probability
weights, and the normalized conditional optimization. A suffix-path
bijection proves equality of the complete continuation problems. Strict
concavity then identifies their unique maximizing policies. Only after
this equivalence is established do we derive the clipped local response,
weighted successor response, and one-sided cap reflection.

Thus the proof no longer moves implicitly from a full history tree to
the public graph. The global knot count, total stored segments, and
fixed-query execution alphabet are separate conclusions with separate
size bounds.

\subsection{3. Closest literature: tree and laminar
allocation}\label{closest-literature-tree-and-laminar-allocation}

Section 2.2 cites and discusses Mjelde (1983) and Tang (1990), rather
than treating chain-nested allocation as the closest class. Proposition
2.1 answers the first question affirmatively: after the positive change
of coordinates \(z_h=w_ha_{v(h)}x_h\), the history formulation is a
separable concave allocation problem with descendant-sum constraints.
The objective coefficients, local bounds, continuation caps, and root
equality are mapped explicitly.

For the referee's questions about memoization, the answer is also
affirmative after normalization. An occurrence response in unconditional
allocation units is its reach weight times the conditional public-vertex
response. Memoizing the normalized subtree responses gives the same
recursion. We do not invent a stochastic obstruction to that
equivalence, and we do not claim to originate scalar marginal allocation
or memoization.

The quantitative results now concern complete responses on a compact
input, coefficient bit heights, and exact execution-state minimization.
We prove those statements directly; they are not justified by asserting
that no earlier paper could contain a related reduction. The global
response-knot bound and the running-maximum label bound follow
transparently from the normalized recursion. Their practical execution
alphabet can be smaller, and Theorem 4.1 characterizes it.

Tang's publisher abstract gives at most twice the number of variables in
one-variable convex subproblems. We report that count as a
scalar-subproblem count, not as a complete rational arithmetic bound. We
do not transfer ascending-chain complexity bounds to arbitrary branching
laminar constraints. Nor do we claim to have established the fastest
published exact algorithm for every tree input. Full-text access to all
closest classical sources was not obtained in this work session;
\texttt{LITERATURE\_AUDIT.md} states this retrieval boundary explicitly.
Consequently, exhaustive historical priority is not marked as
independently verified. The mathematical comparisons and new proofs do
not depend on inventing an unsupported prior complexity bound.

The broader feasible-set connection is also supplied. EC Section EC.1
proves that translating lower bounds gives a laminar polymatroid and
that the root equality selects a base of a truncated rank function.
Federgruen and Groenevelt (1986) and Groenevelt (1991) are discussed
with their discrete/continuous distinctions intact.

\subsection{4. Formal compact-DAG quotient and input
model}\label{formal-compact-dag-quotient-and-input-model}

Theorem 3.1 states that any two occurrences of the same public vertex at
the same incoming scalar price have identical normalized continuation
plans. The proof uses an explicit suffix-history bijection and preserves
probabilities, discounts, local rewards, intervals, and caps. The
ancestor path can affect the price, but cannot affect the conditional
problem after that price is specified.

The statement distinguishes public vertices \(N\), public edges \(M\),
total rational encoding length \(L\), and unfolded occurrences \(H\). It
gives arithmetic work \(O((N+M)N\log(N+1))\), scalar storage
\(O(N^2+M)\), and the bit model in Theorem 3.3. An explicit-tree
implementation must read its \(H\) occurrence records. This is a
comparison of representations, not an impossibility theorem for other
compact-graph methods. A normalized memoized classical implementation
receives the same quotient bound.

\subsection{5. Arithmetic operations versus bit
complexity}\label{arithmetic-operations-versus-bit-complexity}

We have taken the stronger of the referee's two suggested approaches.
Theorem 3.3 supplies a conservative polynomial bit bound under an
explicit binary rational encoding model. Slopes and intercepts are
merged directly as coefficient events. New cap crossings determine
intervals but are not fed back into coefficient interpolation.

On any response interval, every coefficient is a sum of source terms
multiplied by acyclic path probabilities and discounts. A common
denominator divides a fixed product of primitive numerators and
denominators raised to an order-\(N\) power. The normalization bounds
total path weight by the number of visited vertices. This gives bit
height \(S=O(NL+\log(N+M+1))\) and the conservative preprocessing bound
\(O(AS^3)\) for the stated arithmetic count \(A\). This is not a
strongly polynomial claim.

EC Section EC.4 separately bounds execution and aggregate certificate
heights. Summing rewards across distinct prices can create larger
denominators than those in a single response coefficient. Tables 3 and 4
report measured coefficient sizes and separate construction, inversion,
execution, and audit times. The comparison does not treat exact Fraction
arithmetic as constant cost.

\subsection{6. Memory theorem and architecture
accounting}\label{memory-theorem-and-architecture-accounting}

Theorem 4.1 is no longer restricted to the constructed renewal family.
At each public vertex it equates two reached prices precisely when they
induce the same local action and the same behavior classes at every
successor. Backward induction makes this equivalent to equality of the
full future action plan. The smallest common writable alphabet equals
the largest number of distinct behavior classes at any public vertex. A
constructive transition rule attains the bound. The monotonicity of all
future actions also shows that classes are contiguous in the
reached-price order.

An explicit exact example demonstrates strict improvement over counting
numerical prices: two different terminal prices need only one symbol
because their complete behavior is identical. The input graph and
price/threshold/action tables are read-only storage, not free total
implementation storage. Current public vertex and date are observed. Any
retained promise, previous tier, branch identity, or random seed used to
distinguish histories counts as writable state unless explicitly
included in the public state. Root-price recomputation is outside the
fixed-query alphabet result and is measured separately.

\subsection{7. Service-contract
implications}\label{service-contract-implications}

Proposition 4.3 supplies a comparative static tied to continued
participation rather than to a change in mean resource availability. In
the renewal family, hold the public process, branch probabilities, and
mean payment fixed, and spread continuation caps around their common
mean. With a common service-benefit component, this changes
outside-option dispersion at fixed mean outside option. For every
insufficient deterministic alphabet, its optimal value loss is strictly
increasing along this dispersion path.

At exactly zero dispersion one symbol suffices. Any positive dispersion
of distinct branch caps requires distinct terminal actions for exact
optimality, although the value loss tends to zero as dispersion
vanishes. Thus exact implementability and the economic value of
additional retention need not move on the same scale. The statement is
proved for the specified family; we do not present a synthetic
parameterization as a calibrated service application.

\subsection{8. Zero-release shared
table}\label{zero-release-shared-table}

Section 5 keeps the globally shared-table formulation as a supporting
comparator. Proposition 5.1 has public-graph conditional-payment
equations and one decision per information cell. The table is not
reselected independently in successor subproblems. Its convexity is
recognized as standard, rather than advertised as the principal
innovation.

The number of information cells is now \(C_{\mathrm{info}}\), the
restricted value is \(K^{\mathrm{tab}}\), and the minimal execution
alphabet is \(K_*\). The former conflicting use of \(K\) has been
removed from the new manuscript. The longer switching/shared-table
machinery remains available in the retained theory volumes.

\subsection{9. Fixed positive-release oracle versus outer
design}\label{fixed-positive-release-oracle-versus-outer-design}

Proposition 5.2 states an exact value and global supporting inequality
for a fixed table center and release radius. The same finite-price
compiler handles the resulting local intervals. The supporting
inequality is a standard dual-sensitivity construction with the correct
weighted corridor multipliers.

The manuscript does not give the continuous outer shared-table design
problem the preprocessing guarantee of the fixed-query compiler. No
finite exact bound on the number of outer Benders cuts is established or
reported as established. This distinction is explicit in the main text
and response matrix, while the positive fixed-query result remains
intact.

\subsection{10. Algorithmic comparisons and
scale}\label{algorithmic-comparisons-and-scale}

\textbf{10.1. Tree allocation.} \texttt{baselines.py} implements an
independent occurrence-indexed scalar marginal reduction. It constructs
the unfolded laminar input and forms curves by interpolation without
importing the quotient compiler, without memoization, and without
receiving an optimum from the quotient method. All eight comparison
instances agree in exact rational value. This is a faithful
implementation of the stated tree marginal reduction, not a claimed
reproduction of unpublished code or a line-by-line reimplementation of
Mjelde or Tang. A normalized memoized implementation would be
algorithmically equivalent to the quotient recursion.

\textbf{10.2. Generic convex optimization.} The unfolded quadratic
program is solved with SciPy's sparse trust-constr method, analytic
gradient, diagonal Hessian, continuation rows, root equality, and box
bounds. All eight final runs report success. The maximum absolute value
discrepancy is approximately \(3.64\times10^{-7}\); the maximum primal
residual is approximately \(2.62\times10^{-14}\). Actual stopping
settings, iterations, statuses, software versions, and raw repetitions
are included. We do not call a floating-point optimizer an exact
rational certificate. This is an open-source general convex optimization
baseline, not an OSQP or commercial-solver run.

\textbf{10.3. Promise-state recursion.} An independently implemented
primal promise-grid recursion operates on the compact public graph.
Successor promises use grids of resolutions 10, 20, 40, and 80; the
current control absorbs the exact remaining promise. Returned policies
are replayed with rational arithmetic. On six cases, the finest-grid
maximum feasible loss is approximately \(9.77\times10^{-5}\). This
compares discretized primal states with finite exact dual closure
without unfolding the public process. It does not claim that an
arbitrary uniform grid attains exact optimality.

\textbf{10.4. Policy graphs and decomposition.} The preceding baseline
is an actual policy-graph state method. We explain why it is not an SDDP
benchmark: the present output specification is an exact single-query
rational certificate, whereas general sampled decomposition uses
tolerance and sampling choices. There is no relative speed claim about
an unimplemented SDDP method. The published policy-graph literature
remains credited for avoiding naive scenario enumeration.

\textbf{10.5. Larger public graphs.} The scale study goes to 2,045
vertices, rather than stopping at 190. The largest late-renewal instance
has 8,164 edges, 22,419 stored segments, coefficient height 2,213 bits,
and a minimal four-symbol execution alphabet. Its three-run median
construction time is about 2.616 seconds in this environment. A separate
511-vertex breakpoint-dense family stores 197,369 segments; it prevents
the largest graph's small global knot universe from disguising
response-table growth.

Construction, root inversion, certificate generation, audit, machine
minimization, 101-query inversion batches, peak process RSS, and raw
repetitions are separately recorded. Comparison phases have three
repetitions. Scale certificate/audit phases have one explicitly labeled
diagnostic run. Censored pilot attempts are retained outside the
completed timing tables. Peak RSS includes imports and is not mislabeled
as incremental algorithm storage. These measurements establish a
specified comparison, not universal dominance or a dedicated-hardware
performance claim.

\subsection{11. Exponential history
size}\label{exponential-history-size}

The main narrative no longer uses the number of histories as a surrogate
for outperforming modern stochastic programming. Section 2 explains that
both public-graph state methods and normalized tree memoization exploit
recombination. The residual exactness statement is the finite response
closure, the binary-input construction, and the minimal execution
machine. The experimental comparator includes a primal state method
using the same compact public graph.

\subsection{12. Focus without deleting historical
results}\label{focus-without-deleting-historical-results}

The new main manuscript is 22 PDF pages including its title page,
references, and four tables; the new electronic companion is 10 pages.
The principal proofs appear in the main manuscript. Optional
implementation and boundary details appear in the companion.

No historical derivation was removed to achieve this focus. The R29 main
source/PDF and companion source/PDF have byte-identical copies in
\texttt{predecessor/}; all R1--R29 source and result directories remain.
The complete general switching, sensitivity, release, continuous-state,
cache, and governance results can therefore still be reviewed, not
merely read as a prose summary. \texttt{PRESERVATION\_REPORT.json}
verifies every one of the 1,540 inherited files, allowing changes only
to six reader entry points whose originals are retained.

\subsection{13. Technical details of the finite-price
theorem}\label{technical-details-of-the-finite-price-theorem}

\textbf{13.1. Capped versus pre-cap values.} The definitions now
distinguish \(H_v\), \(H_v^{\mathrm{pre}}\), \(R_v\), and \(d_v\). Only
the current cap is removed in the pre-cap problem; successor caps stay
imposed.

\textbf{13.2. Root equality and cap.} The root has its own cap, and the
theorem requires \(b_0\in[L_0,U_0]\subseteq(-\infty,B_0]\). Feasible
payment intervals are computed explicitly before the root query.

\textbf{13.3. Plateaus.} The leftmost finite crossing is justified,
including a cap at the minimum payment. If two prices produce the same
payment, the two optimality inequalities and strict concavity imply the
same full optimizer. Therefore dual nonuniqueness does not produce
ambiguous actions.

\textbf{13.4. Occurrence invariance.} This is the opening claim and
proof of Theorem 3.1, not an unstated induction premise.

\textbf{13.5. Knot universe versus stored lines.} The global union is at
most \(3N\), while storing a separate response at every vertex can
require \(O(N^2)\) scalar records. Both bounds remain visible in theory
and tables.

\textbf{13.6. Multiple coordinates.} Corollary 3.2 gives at most
\(2D+N\) response knots but at most \(N+1\) fixed-query price labels.
Local clipping knots describe an action response; only a cap barrier
changes the stored running price. These are different counts.

The informal \(d_v(-\infty)\) notation is replaced by the explicit
finite left-tail payment \(d_v^-\).

\subsection{14. Memory-frontier details}\label{memory-frontier-details}

\textbf{14.1. Controller observations.} Section 4.1 specifies the public
vertex/date, read-only data, writable register, and transition
observations. Carrying a promise, predecessor identity, inherited
action, or persistent random seed is not allowed to evade the alphabet
count by renaming it. The renewal frontier additionally specifies the
common terminal observation and absence of a free branch-revealing
physical signal.

\textbf{14.2. Randomization.} Theorem 4.1 extends the necessary alphabet
for exact attainment of the full optimum. Conditional averaging
preserves the linear feasibility requirements, and strict concavity
forces a randomized exact-optimal controller to use the unique action at
every positive-probability public history almost surely. Distinct
continuation behaviors therefore need disjoint support in writable
symbols. Retained randomness is counted. This argument does not
establish equality of every randomized and deterministic
restricted-alphabet value; Theorem 4.2 states its frontier for
deterministic encodings.

\textbf{14.3. Broader structure.} Theorem 4.1 applies the behavioral
quotient throughout the finite-price class. Theorem 4.2 independently
extends the ordered renewal frontier to unequal positive branch
probabilities, strictly increasing concave terminal rewards, and
heterogeneous convex nondecreasing intermediate costs. Exact enumeration
compares 2,308 arbitrary, not necessarily contiguous, partitions against
the ordered recurrence. The original eight-branch frontier is recovered
exactly as a special case.

\subsection{15. Normal cones at corridor ties and
singletons}\label{normal-cones-at-corridor-ties-and-singletons}

The proof of Proposition 5.2 and EC Section EC.7 give the complete
split. A lower interval normal can be assigned among original lower
inequalities attaining the effective endpoint, and an upper normal among
original upper inequalities attaining its endpoint. At a singleton both
orientations are available. Any nonnegative split preserving the total
normal preserves stationarity and complementarity.

The global supporting inequality follows by applying the old multipliers
to arbitrary policies feasible under new parameters; it does not presume
those policies use the old price machine. Exact tests cover 150 feasible
anchors, including fixed tiers, coincident physical/corridor bounds, and
zero release width. They contain 96 tied-normal occurrences and 1,185
exact supporting inequalities under physical-first, corridor-first, and
equal splits. The statement remains a fixed-query supergradient result,
not a finite-iteration outer-design theorem.

\subsection{16. Literature and
positioning}\label{literature-and-positioning}

The main references and literature audit now include Mjelde (1983), Tang
(1990), Federgruen and Groenevelt (1986), and Groenevelt (1991),
followed by the nested-allocation and policy-graph lineages. Proposition
2.1 and EC Section EC.1 map the actual feasible set, rather than compare
labels alone. The continuous quadratic, discrete allocation, explicit
tree, compact DAG, and general policy-graph input models are not
conflated.

The audit records exact bibliographic identifiers and which publisher
text was accessible. It does not purport to exclude all possible
historical antecedents. Our affirmative case is the stated package of
proved compact-response, bit-complexity, minimal-memory, and dispersion
results, not an unsupported claim that classical marginal reasoning was
absent.

\subsection{17. Requested ingredients for a compelling
resubmission}\label{requested-ingredients-for-a-compelling-resubmission}

\textbf{A. Priority and laminar allocation:} exact mapping, normalized
memoization equivalence, explicit compact-input results, and broader
rank formulation are provided. Exhaustive priority against all full
texts is not claimed verified.

\textbf{B. Complexity:} a rational bit theorem replaces reliance on an
arithmetic count alone; certificate heights and measured bit growth are
separate.

\textbf{C. Baselines:} independent tree reduction, sparse general convex
optimization, and a public-graph promise-state method are implemented
and reported. SDDP is discussed without invented timings or dominance
claims.

\textbf{D. Memory:} the exact minimal behavioral machine is
characterized for the finite-price class; the ordered frontier is
generalized beyond the original equal-probability quadratic family.

\textbf{E. Focus:} a 22-page main manuscript and 10-page companion
concentrate on the new center, while all broader derivations remain
accessible in full.

\textbf{F. Operational implication:} fixed-mean outside-option
dispersion strictly raises the value loss from insufficient retention in
the stated renewal family, with an exact-alphabet discontinuity at zero
dispersion.

\subsection{18. Presentation and remaining minor
points}\label{presentation-and-remaining-minor-points}

The title refers to compact allocation quotients and minimal price
memory. The abstract says that the finite alphabet belongs to a
fixed-query implementation; it does not imply a finite price set for all
possible root promises. The abstract is text only and has 191 words. The
introduction has no equations or mathematical notation. The manuscript
uses author--year citations, an alphabetical bibliography, 11-point
text, one-and-a-half spacing, one-inch margins, and no footnotes. The
submission checklist distinguishes formatting checks from author
declarations that cannot be inferred from repository files.

The prior adverse cache result remains explicit: 83 of 84
strict-tolerance queries refreshed, without a measured speed advantage.
The exact verifiers are described as implementation and certificate
checks, not as independent mathematical proofs. Source-preservation
logistics appear in the evidence documents rather than replacing
scientific argument.

\subsection{19. Concluding response to the
editor}\label{concluding-response-to-the-editor}

The revised paper accepts the correct classical allocation connection
and proves the quantitative and operational statements on which its
contribution now rests. In addition to the R29 finite response
construction, it supplies an explicit quotient proof, polynomial
rational preprocessing, a minimal execution-alphabet theorem, a
heterogeneous memory frontier, a fixed-mean dispersion comparative
static, and a complete tied-bound supporting-cut argument. The
computational study now contains actual algorithmic alternatives and
distinct scale families with auditable limitations.

We submit these results for renewed mathematical and editorial scrutiny.
We do not equate a passing code audit with publication readiness, nor
mark an unperformed external literature verification as complete. The
complete prior theory and evidence are retained so that the next referee
can evaluate the scientific changes directly against the reviewed
baseline.

\end{document}
